\setcounter{section}{1}
\setlength{\parindent}{0pt}  % niente rientro
\setlength{\parskip}{6pt}
\section{Requisiti funzionali}
\subsection{Requisiti comuni}
\subsubsection*{RF-01-Registrazione utente}
Il sistema deve consentire agli utenti di potersi registrare autonomamente direttamente dalla webapp. L'utente accede ad una sezione dedicata alla registrazione. Il sistema deve richiedere i dati necessari:
\begin{itemize}
    \item Nome
    \item Cognome
    \item Email
    \item Codice fiscale
    \item Numero di telefono
    \item Comune di residenza
    \item Password
\end{itemize}
Il sistema deve validare i dati inseriti dall'utente, verificando che:
\begin{itemize}
    \item Tutti i campi obbligatori siano stati compilati
    \item Il nome e il cognome non contengano numeri o caratteri speciali
    \item La email sia in un formato valido e che non sia già registrata nel sistema
    \item Il codice fiscale sia in un formato valido e che non sia già registrato nel sistema
    \item Il numero di telefono sia in un formato valido
    \item La password rispetti i criteri di sicurezza (lunghezza minima, presenza di caratteri speciali, ecc.)
\end{itemize}
Conclusa la validazione, vengono creati i profili utente e inviata una mail di conferma all'indirizzo fornito.
%TODO: aggiungere requisito per accettazione termini e condizioni/privacy policy?
%TODO: aggiungere requisito per captcha?
%TODO: aggiungere verifica del numero di telefono tramite sms con codice OTP?
\subsubsection*{RF-02-Login utente}
Il sistema deve consentire agli utenti di effettuare il login alla webapp utilizzando le credenziali fornite in fase di registrazione (email e password) o tramite autenticazione OAuth2 con provider esterni (Google, Microsoft) se collegato. Il sistema deve verificare le credenziali inserite e, in caso di successo, reindirizzare l'utente alla sua homepage personalizzata. In caso di fallimento, il sistema deve mostrare un messaggio di errore appropriato.
%TODO: verificare se sia meglio dire "recupero password" o "cambio password", o creare due requisiti distinti, il cambio password potrebbe essere fatto anche da utente loggato dal proprio profilo utente, mentre recupero password non e' un vero e proprio recupero ma un cambio password tramite link inviato via email
\subsubsection*{RF-02.1-Recupero password}
Il sistema deve fornire una funzionalità di recupero password per gli utenti che hanno dimenticato le loro credenziali di accesso. L'utente deve poter richiedere il reset della password inserendo il proprio indirizzo email. Il sistema deve inviare una email con un link sicuro per reimpostare la password. Una volta cliccato il link, l'utente deve essere reindirizzato ad una pagina dove può inserire una nuova password, che deve rispettare i criteri di sicurezza definiti nel requisito RF-01.
\subsubsection*{RF-03-Accesso all'area personale}
Il sistema deve consentire agli utenti autenticati di accedere ad un'area personale dedicata, dove possono visualizzare le informazioni del proprio profilo e gestire le impostazioni dell'account. L'area personale deve essere protetta e accessibile solo agli utenti autenticati.
\subsubsection*{RF-03.1-Modifica dati personali}
Il sistema deve consentire agli utenti di poter modificare o eliminare, a patto che non siano dati obbligatori, i propri dati personali, tramite l'area personale dedicata come specificato da requisito RF-03.
%NON FATTO\subsubsection*{RF-03.2-Modifica email}
%Il sistema deve consentire agli utenti di modificare l'indirizzo email associato al loro account, tramite l'area personale dedicata come specificato da requisito RF-03. Il sistema deve inviare una email di conferma al nuovo indirizzo per verificare la validità dello stesso. Una volta confermato, il sistema deve aggiornare l'indirizzo email nel profilo dell'utente.
%TODO: prevedere inoltre una verifica del nuovo numero di telefono tramite sms con codice OTP?
%NON FATTO\subsubsection*{RF-03.2-Modifica numero di telefono}
%Il sistema deve consentire agli utenti di modificare il numero di telefono associato al loro account, tramite l'area personale dedicata come specificato da requisito RF-03. Il sistema deve validare il nuovo numero di telefono per assicurarsi che sia in un formato corretto prima di aggiornare le informazioni dell'utente.
\subsubsection*{RF-03.2-Cancellazione dell'account}
Il sistema deve consentire agli utenti di poter cancellare il proprio account in modo permanente, tramite l'area personale dedicata come specificato da requisito RF-03. Prima di procedere con la cancellazione, il sistema deve richiedere una conferma esplicita dall'utente per evitare cancellazioni accidentali. Una volta confermata, il sistema deve eliminare tutti i dati associati all'account dell'utente in conformità con le normative sulla protezione dei dati.
\subsubsection*{RF-04-Logout utente}
Il sistema deve consentire agli utenti autenticati di effettuare il logout dalla webapp in modo sicuro. Una volta effettuato il logout, l'utente deve essere reindirizzato alla pagina di login o alla homepage pubblica del sito e non deve piu' avere accesso alle funzionalità riservate agli utenti autenticati senza effettuare nuovamente il login.
\subsubsection*{RF-05-Visualizzazione dati energetici personali}
Il sistema deve consentire all'utente di visualizzare i propri dati energetici personali, tra cui consumi, stato batteria degli impianti, risparmio co2 e produzione nel tempo. I dati devono essere presentati tramite grafici e valori numerici, con aggiornamento periodico o in tempo reale se disponibile.
\subsubsection*{RF-06-Reportistica utente}
Il sistema deve consentire all'utente di scaricare un report in formato pdf della propria produzione storica o in tempo reale.

\subsubsection*{RF-07-Gestione impianti energetici}
Il sistema deve consentire all'utente di registrare nuovi impianti produttivi di energia rinnovabile (es. fotovoltaico, eolico). Ogni impianto deve avere almeno: un codice identificativo, potenza nominale, posizione geografica, stato operativo e data di installazione. L'utente deve poter sospendere temporaneamente un impianto per manutenzione o rimuovere un impianto non piu' operativo.

\subsubsection*{RF-08-Ricerca CER}
Il sistema deve permettere ad un utente di cercare una CER attraverso il suo nome o codice identificativo, visualizzando un anteprima della CER ricercata.

\subsubsection*{RF-09-Richiesta adesione CER}
Il sistema deve permettere ad un utente di fare richiesta di adesione per una CER, verrà posto nella lista delle richieste d'adesione della CER, in attesa di una risposta da parte di un amministratore di quella comunità.

\subsection{Requisiti per Admin della CER}
\subsubsection*{RF-10-Gestione membri della CER}
Il sistema deve consentire all’amministratore della comunità di visualizzare l’elenco dei membri appartenenti alla CER, aggiungere nuovi membri e modificare o rimuovere quelli esistenti. Per ogni membro devono essere mostrati i dati identificativi, i punti di consumo o produzione associati e lo stato della loro partecipazione (attivo, in attesa, sospeso). Le operazioni di modifica e rimozione devono essere tracciate nel registro delle attività.
\subsubsection*{RF-11-Gestione richieste di adesione}
Il sistema deve consentire all’amministratore di approvare o rifiutare le richieste di adesione alla CER. Le richieste devono contenere i dati dell’utente e le informazioni necessarie alla valutazione (es. nome, cognome, indirizzo dell’abitazione). In caso di approvazione, il sistema deve inviare una notifica email all’utente e aggiornarne lo stato a “membro attivo”.
\subsubsection*{RF-12-Visualizzazione dati energetici e reportistica}
Il sistema deve consentire all’amministratore di visualizzare report aggregati sulla produzione e il consumo energetico della CER. L’interfaccia deve permettere la selezione di un intervallo temporale (giornaliero, mensile, annuale) e la visualizzazione di grafici e tabelle riassuntive. I dati devono poter essere esportati in formato PDF.
\subsubsection*{RF-13-Gestione bacheca avvisi}
Il sistema deve consentire all'amministratore di gestire la bacheca avvisi della comunità, eliminando o aggiungendo nuovi avvisi per i membri.
\subsubsection*{RF-14-Visualizzazione notifiche}
Il sistema deve permettere all'amministratore di visualizzare i messaggi inviati dai membri della sua CER in un area dedicata.







%FUNZIONALITA' COMUNE NON PRESENTI
%\subsection{Requisiti per il Rappresentante del Comune}
%\subsubsection*{RF-11-Gestione e supervisione delle CER comunali}
%Il sistema deve consentire al rappresentante del Comune di visualizzare tutte le CER attive nel territorio comunale, accedendo a informazioni sintetiche come numero di membri, produzione complessiva, energia condivisa e stato operativo. Il rappresentante deve poter consultare i dettagli di ciascuna CER ma non modificarne la configurazione interna.
%\subsubsection*{RF-12-Visualizzazione dati aggregati}
%Il sistema deve fornire al rappresentante del Comune una dashboard riepilogativa con i principali indicatori energetici aggregati sul territorio comunale: energia prodotta, energia condivisa, numero di cittadini coinvolti e riduzione totale di emissioni. I dati devono poter essere filtrati per periodo o per singola CER.
%\subsubsection*{RF-13-Gestione richieste di approvazione o supporto}
%Il sistema deve consentire al rappresentante del Comune di visualizzare le richieste provenienti dalle CER per attività che necessitano di supervisione (es. nuove comunità da validare, segnalazioni tecniche o amministrative). Il rappresentante deve poter aggiornare lo stato della richiesta e lasciare un commento visibile agli amministratori della CER.
%\subsubsection*{RF-14-Comunicazioni istituzionali}
%Il sistema deve consentire al rappresentante del Comune di inviare comunicazioni o avvisi a tutte le CER registrate nel territorio. Le comunicazioni devono poter includere testo e allegati (PDF, immagini) e devono essere visibili nell’area notifiche delle rispettive comunità.
\subsection{Requisiti per i membri della CER}
\subsubsection*{RF-15-Notifiche e avvisi}
Il sistema deve notificare ai membri eventi rilevanti, come l’aggiunta di una comunicazione, un avviso o l’approvazione di una richiesta. I membri vengono notificati via posta elettronica.
\subsubsection*{RF-16-Partecipazione e segnalazioni}
Il sistema deve consentire ai membri di inviare segnalazioni o richieste di supporto tecnico/amministrativo all’amministratore della CER via messaggio.
\subsubsection*{RF-17-Abbandona CER}
Il sistema deve consentire al membro di una CER di abbandonarla in qualsiasi momento. 

\subsection{Requisiti per l'Amministratore di Sistema (Staff Rinova)}

\subsubsection*{RF-AS-01-Registrazione Nuova CER}
Il sistema deve consentire all'amministratore di sistema di registrare una nuova Comunità Energetica Rinnovabile sulla piattaforma. L'operazione richiede l'inserimento dei dati anagrafici obbligatori (denominazione, codice fiscale/P.IVA, indirizzo legale, area geografica di riferimento) e la definizione dei parametri tecnici iniziali. Contestualmente alla creazione, il sistema deve permettere di associare o invitare un utente che assumerà il ruolo di "Amministratore della CER".

\subsubsection*{RF-AS-02-Modifica Dati Anagrafici CER}
Il sistema deve consentire all'amministratore di sistema di aggiornare le informazioni di una CER esistente. Questo include la modifica dei dati anagrafici, dei contatti di riferimento e la possibilità di riassegnare il ruolo di Amministratore della CER a un altro utente in caso di avvicendamento gestionale interno alla comunità.

\subsubsection*{RF-AS-03-Eliminazione CER}
Il sistema deve consentire all'amministratore di sistema di eliminare definitivamente una CER dalla piattaforma. L'eliminazione deve comportare la rimozione o l'anonimizzazione dei dati sensibili associati alla comunità, in conformità con le policy di privacy.

\subsubsection*{RF-AS-04-Validazione impianto}
Il sistema deve consentire all'amministratore di sistema di validare l'aggiunta di un impianto, modificandone lo stato in "attivo".